\scanpage{171}
\mspage{484}
% position: upper paragraph
(${}^{110}$) (11 octobre) J'ai eu envie depuis un jour ou deux de faire le point, en quelques mots, où en est (après huit ans) cette ``restauration progressive de l'équilibre yin-yang'' en moi.

% position: middle long paragraph
Peut-être le changement le plus important de tous est dans une \mbox{\unsure{a c c e p|a c c e p-}}\linebreak\mbox{t a t i o n} beaucoup plus grande que par le passé de ma personne telle qu'elle est vraiment d'instant en instant. Une autre façon de l'exprimer, c'est que les mécanismes de répression en moi se sont considérablement assouplis. Comme je \unsure{l'a|l'ai} dit hier, certains ont disparu après avoir été découverts et compris, et d'au-\linebreak tres, que j'avais ignorés ma vie durant, me sont devenus familiers dans leurs manifestations de tous les jours. Je les vois en action, non comme des ennemis qu'il me faudrait essayer d'extirper coûte que coûte, mais comme faisant partie de la multiplicité des facettes de mon être conditionné, et par là, de la \unsure{riches|riches-}\linebreak se du ``donné'' présent, lequel reflète fidèlement mon histoire passée; aussi bien l'histoire ``ancienne'' de mes conditionnements et des racines de la division dans mon être, que l'histoire plus récente de ma maturation, du travail donc par le-\linebreak quel je finis par déballer et par ``manger'' et assimiler le paquet initial légué par mes parents et par leurs successeurs. Cette ``acceptation'' en moi inclut donc\unsure{:} non seulement des pulsions et traits de ``l'enfant'' que j'avais pendant longtemps ignorés et réprimés (et notamment ceux qui reflètent les aspects féminins en moi), mais également les mécanismes de répression propres au ``patron'', c'est-à-dire justement des mécanismes invétérés de ``non-acceptation''! Accepter ces der-\linebreak niers n'a rien de commun avec ``les cultiver'', ou les fortifier. Au contraire, c'est un premier pas indispensable pour les dénouer ou les désamorcer tant soit peu, par l'effet d'une attention curieuse et aimante. L'expérience de ces huit années me donne la conviction que, pour peu que cette attention-là plonge assez profond et jusqu'à la racine même de la répression, celle-ci se résoud et dis-\linebreak paraît en libérant une énergie considérable - celle qui jusque là était immobili-\linebreak sée pour maintenir contre vents et marées tel ensemble de mécanismes répressifs, et les habitudes de pensée et autres qui servent à les maintenir.

% position: lower paragraph
Mais ce n'est pas vis-à-vis des aspects par nature ``noués'' de ma personne, que cette acceptation nouvelle de moi-même a d'abord fait son apparition dans ma vie. Elle est venue sans tambour ni trompette, dès avant la découverte de la méditation, donc dès avant les ``retrouvailles'' la suivant de près. C'était au mois de juillet 1976, au cours d'une courte liaison amoureuse avec une jeune femme, G. , peut-être un brin plus ``homasse'' dans ses façons d'être que les femmes que j'avais aimées précédemment. Le hasard (?) a voulu que les circons-\linebreak tances matérielles qui ont entouré ces amours étaient telles, que je me voyais placé dans un rôle typiquement ``féminin''. Je faisais le ménage et préparais les
% position: foot left short stroke
\margin{--}
% position: foot pencil
\margin{170}
